\par
\SourceLeaf{488}%
Ich habe noch hinzuzufügen, dass ich über diesen Gegenstand bereits vor einem Jahre Sr. — — — ein ausführliches Memoire vorgelegt habe, und dass dieser, die Wichtigkeit der Fortsetzung der dänischen Messungen nicht verkennend, wenigstens sofort die Möglichkeit derselben sicherte, indem er mir den Befehl ertheilte, Lüneburg an die zum Theil in vergänglichen Signalpunkten bestehenden südlichen Punkte jener Messung anzuschliessen. Dies ist im vorigen Herbst geschehen, und dadurch die künftige Fortsetzung von dem Untergang der precären Punkte unabhängig gemacht.

Wenn Ew. — — diesen Ideen Ihren Beifall schenken, und sie würdig halten, für ihre Realisirung zu wirken, so muss ich schon die Beabsichtigung einer möglichst vollkommenen und neben jeder andern ehrenvoll bestehenden Ausführung voraussetzen, und in dieser Hinsicht würde ich mich der Ausführung, wenn sie mir anvertrauet würde, mit Vergnügen unterziehen, und eine vielleicht länger als Einen Sommer dauernde Unterbrechung meiner rein astronomischen Arbeiten für kein Opfer halten.

Was den zur Verlängerung der dänischen Gradmessung durch das Königreich Hannover erforderlichen Kostenaufwand betrifft, so ist es wegen seiner Abhängigkeit von mancherlei nicht vorherzusehenden Umständen freilich unmöglich, denselben mit einiger Genauigkeit im Voraus anzugeben. So hat die Witterung auf die Zeitdauer, und dadurch auf die Kosten einen wesentlichen Einfluss. Ich wüsste nicht, dass von irgend einer Gradmessung die Kosten öffentlich bekannt gemacht wären, und wenn ich gleich von dem Dänischen Astronomen über die bisherigen Kosten seiner Gradmessung Mittheilung erhalten könnte, so würden sich doch daraus die Kosten ähnlicher Operationen im Königreich Hannover, wegen der grossen Verschiedenheit der Localverhältnisse, nur sehr unsicher schätzen lassen. So werden z. B. im Hannöverschen die Kosten für Fuhren Behuf Transports der Personen und Instrumente einen sehr bedeutenden Theil der Gesammtkosten austragen, da jene im Dänischen, wo die Fuhren ex officio in natura geleistet werden, fast ganz aus der Rechnung wegfallen. Doch glaube ich, alles wohl erwogen, dass die Summe von 1500 Pfund Sterling zur Bestreitung aller Kosten hinreichen würde, und versteht sich von selbst, dass darüber demnächst Rechnung abgelegt werden würde.

Der Professor \textsc{Schumacher} hat bei seiner Gradmessung ausser ein Paar Amanuensen, zwei Officiere von Capitains Rang zu Gehülfen, deren Geschäft es ist, die Gegend vorher zu bereisen, schickliche Punkte für Dreiecksstationen auszusuchen, um sie, nach vorläufig daselbst gemachten Messungen dem Prof. S. zur Auswahl vorzuschlagen, hernach auf den ausgewählten Punkten die Errichtung von Signalen, wo es nöthig ist, und andere nöthige Vorkehrungen zu betreiben, mit einem Wort, alle Vorbereitungen zu machen, dass der Astronom überall ohne vielen Zeitverlust zum Beobachten schreiten kann; endlich da, wo es nöthig ist, bei den Beobachtungen die Nebenoperationen zu übernehmen. Die erforderlichen Eigenschaften für solche Gehülfen sind daher, nicht sowohl besonders tiefe mathematische oder astronomische Einsichten, als vielmehr reger Eifer für die Sache, die grösste Pünktlichkeit und Sinn für die grösste Genauigkeit, eine gewisse praktische Anstelligkeit, einige Kenntnisse vom Bauwesen, einige Bekanntschaft mit dem Geschäftsgange in unserm Lande bei denjenigen Behörden, mit welchen in solchen Angelegenheiten Berührungen vorkommen. Ein grösseres Personal als bei der Dänischen Gradmessung würde auch bei der Hannöverschen nicht nöthig sein, und ich würde mir sogleich, sobald eine Resolution gefällt ist, angelegen sein lassen, taugliche Gehülfen selbst auszusuchen. Natürlich würde es aber von Ew. — — Ermessen abhängen, ob vielleicht rathsam erachtet würde, in \textit{der oben erwähnten} Beziehung noch mehrere Personen als Volontairs den Messungen beiwohnen zu lassen, um sich zu feinern geodätischen Arbeiten vorzubereiten. — —

Göttingen den 30. Mai 1819.
\begin{flushright}\textsc{Carl Friedrich Gauss}.\end{flushright}
